\answer{ex:04:1}{(a)~$\approx1.1\cdot10^{11}$ bits/J. (b)~$\log_2\sqrt{10}\approx1.7$ bit contre $81$: $\approx50$ fois moins.}
\answer{ex:04:2}{$r^*=(P_0/E_{\text{imp}})\ln\bigl(1/(r^*\Delta t)\bigr)$, $P_0/E_{\text{imp}}=1.16\,\text{s}^{-1}$: $r^*\approx6$ impulsions par seconde, $7.4\cdot10^{10}$ bits/J.}
\answer{ex:04:3}{$\mathcal I=\frac{N}{1+(N-1)c}=9.9$ (limite $10$) contre $\mathcal I=\frac{N}{1-c}=1111$: un rapport de $\approx110$; la corrélation n'est nuisible que si elle ressemble au signal.}
\answer{ex:04:4}{$\theta\approx68.7$ et $19.9$ (en unités de $w$), $m=19$ et $m=10$: le destinataire rapide a besoin de presque deux fois moins d'entrées coïncidentes, bien que son entrée moyenne soit cinq fois plus petite.}
\answer{ex:04:5}{$U\tau_{\text{rec}}\ge0.725\,\text{s}$ et $\tau_{\text{rec}}(1-2U)\le0.05\,\text{s}$, ce qui exige $U\ge0.483$; le plus petit $\tau_{\text{rec}}=0.725\,\text{s}$ pour $U=1$ ($1.45\,\text{s}$ pour $U=0.5$).}
\answer{ex:04:6}{(a)~$\theta\,e/(e-1)\approx1.58\,\theta$ quels que soient $\tau$ et $\theta$. (b)~$14$ impulsions.}
\answer{ex:04:7}{(a)~$1.62$ bit par mot: $0.77$ pour le nombre d'impulsions, $0.84$ pour leur disposition. (b)~$1.13$ et $1.59$ bit: sur $30$ mots, l'estimation perd près d'un tiers.}
